\begin{explanationbox}
\textbf{\textcolor{CExplanation}{从你容易看见的竖直差开始}}

先看 $l:y=kx+b$。与 $P(x_0,y_0)$ 具有相同横坐标、又在 $l$ 上的点是
\[
 R(x_0,kx_0+b),\qquad \Delta=y_0-(kx_0+b).
\]
\begin{center}
\begin{minipage}{0.98\linewidth}
\begin{tcolorbox}[colback=white,colframe=CStatement!25,boxrule=0.4pt,
  arc=1mm,left=3mm,right=3mm,top=3mm,bottom=3mm]
\centering
\textbf{同横坐标的点 R，不一定是垂足 H}\par\medskip
\begin{tikzpicture}[font=\small,>=Stealth]

\begin{scope}[x=1.4cm,y=1.4cm]
\coordinate (P) at (0,3);\coordinate (R) at (0,0);\coordinate (H) at (1.44,1.08);
\coordinate (E) at (1.05,0);
\fill[CStatement!6] (R)--(P)--(H)--cycle;
\draw[CStatement,thick] (-.7,-.525)--(3.6,2.7) node[above left] {$l:y=kx+b$};
\draw[CProof,very thick] (R)--(P);
\draw[CExample,very thick] (P)--(H);
\draw[gray,dashed] (R)--(E);
\draw[CExample] (1.296,1.272)--(1.104,1.128)--(1.248,.936);
\pic[draw,angle radius=7mm,angle eccentricity=1.45,"$\theta$"] {angle=E--R--H};
\pic[draw,angle radius=7mm,angle eccentricity=1.45,"$\theta$"] {angle=R--P--H};
\fill (P) circle (2pt) (R) circle (2pt) (H) circle (2pt);
\node[above left] at (P) {$P(x_0,y_0)$};
\node[below right,align=left] at (R) {$R$\\$(x_0,kx_0+b)$};
\node[right] at (1.6,.85) {$H$};
\node[left,text=CProof] at (-.1,1.55) {$PR=\Delta$};
\node[above right,text=CExample] at (.72,2.04) {$PH$};
\end{scope}
\end{tikzpicture}\par\medskip
图示 $k>0$、$\Delta=y_0-kx_0-b>0$。$PR$ 是竖直线段，$PH$ 才垂直于 $l$；在直角三角形 $PHR$ 中，$PH=PR\cos\theta$。
\end{tcolorbox}
\end{minipage}
\end{center}


先取图中的情形 $k>0,\Delta>0$。沿竖直方向从 $R$ 到 $P$ 的长度是 $PR=\Delta$，但点到直线的距离应取\textbf{垂直于直线的线段 $PH$}。
因为 $PH\perp l$，图中两个标为 $\theta$ 的角相等，所以
\[
 PH=PR\cos\theta.
\]
\begin{center}
\begin{minipage}{0.98\linewidth}
\begin{tcolorbox}[colback=white,colframe=CStatement!25,boxrule=0.4pt,
  arc=1mm,left=3mm,right=3mm,top=3mm,bottom=3mm]
\centering
\textbf{根号来自斜率方向三角形}\par\medskip
\begin{tikzpicture}[font=\small,>=Stealth]

\coordinate (U) at (0,0);\coordinate (V) at (4,0);\coordinate (W) at (4,3);
\draw[CStatement,thick] (U)--(V)--(W)--cycle;
\draw (3.75,0)--(3.75,.25)--(4,.25);
\pic[draw,angle radius=10mm,angle eccentricity=1.4,"$\theta$"] {angle=V--U--W};
\node[below] at (2,0) {$1$};
\node[right] at (4,1.5) {$|k|$};
\node[above left,rotate=36.87] at (2,1.5) {$\sqrt{1+k^2}$};
\node[align=left,anchor=west] at (5.25,1.5) {$\displaystyle\cos\theta=\frac{1}{\sqrt{1+k^2}}$};

\end{tikzpicture}\par\medskip
图按 $|k|=3/4$ 等比例绘制。把水平边长记为 $1$，竖直边长就是 $|k|$；由勾股定理，斜边长为 $\sqrt{1+k^2}$。
\end{tcolorbox}
\end{minipage}
\end{center}


因此，竖直差需要乘上 $1/\sqrt{1+k^2}$，才变成垂直距离：
\[
 \boxed{PH=\frac{|y_0-kx_0-b|}{\sqrt{1+k^2}}}.
\]
绝对值使两侧的点都得到非负长度。若 $k<0$，取直线与水平方向的锐角，仍有同一个余弦系数；若 $k=0$，竖直方向恰好垂直于直线，$H=R$，分母为 $1$；若 $\Delta=0$，则 $P=H=R$，距离为零。

\medskip
\textbf{\textcolor{CExplanation}{两个根号各自负责什么？}}

垂直于 $y=kx+b$ 的一个方向向量为 $(k,-1)$，长度是 $\sqrt{1+k^2}$。若把从 $P$ 到 $H$ 的位移写成 $s(k,-1)$，则
\[
 |s|=\frac{PH}{\sqrt{1+k^2}}
     =\frac{|\Delta|}{1+k^2}.
\]
\textbf{第一个根号}把竖直差换算为垂直距离；\textbf{第二个根号}把实际长度换算成向量 $(k,-1)$ 的倍数。$s$ 的正负由点在直线的哪一侧决定，证明三会给出带符号的完整坐标式。
\end{explanationbox}
